\section{Introduction}\label{sec:introduction}


Growing attention to decarbonization and Arctic accessibility has renewed interest in alternative Europe--Asia shipping corridors. The Northern Sea Route (NSR) is frequently presented as a 30--40\% shorter alternative to Suez based on great-circle (GC) geometry \citep{liu2010_nsr,furuichi2015_platform,schoyen2011_nsrecon}. Recent corridor assessments continue to report NSR advantages primarily from geometric or stylized baselines before adding cost or feasibility layers \citep{li2023_transarctic_jmse,zeng2020_competitiveness}. Some studies propagate these geometric savings to time, fuel, or CO$_2$ without first demonstrating sea-only navigability \citep{iaph2013_portplanning,wan2018_sustainability}, whereas others explicitly model Arctic feasibility and operating speeds, often jointly with regulatory or emissions-control layers \citep{SmithStephenson2013,melia2016_future,kavirathna2023_nsr_feasibility}. Because corridor baselines are increasingly used to parameterize network, climate, and policy models, small routing-representation biases can scale into large system-level conclusions \citep{poo2024_arctic_network}. What remains under-examined is how enforcing sea-only feasibility reshapes corridor baselines and their downstream translation to time, fuel, and CO$_2$ under corridor-specific speeds.

We compare geometric GC baselines (direct GC and GC via corridor macro-waypoints) with physically feasible sea-only routes computed on a coastline-masked 0.5$^\circ$ A* graph. The framework is applied to Rotterdam--Yokohama, a representative Europe--Asia pair spanning the Suez, Cape of Good Hope, and NSR corridors. Each corridor is represented by three waypoint variants to capture distinct routing philosophies: (i) \emph{service-guided} (hub/chokepoint oriented), (ii) \emph{bluewater-oriented} (open-ocean favouring), and (iii) \emph{channel/coast-guided} (lane/coast following). This design magnifies routing-method effects over an intercontinental distance while remaining comparable to prior corridor studies \citep{li2023_transarctic_jmse,kavirathna2023_nsr_feasibility,meza2023_arctic_routes}. Distances are mapped to voyage time using corridor-typical service speeds and to fuel/CO$_2$ using explicit main- and auxiliary-engine accounting. To isolate routing effects, we exclude sea ice, weather, emission control areas (ECAs), bathymetry and draft limits, and AIS-based tracking; these layers are reserved for subsequent constrained analyses so that environmental and regulatory factors do not confound the methodological comparison.

Our analysis yields three headline insights. First, enforcing sea-only feasibility preserves the distance ranking (NSR $<$ Suez $<$ Cape): NSR remains shortest once routes are rendered navigable. Second, distance does not translate linearly to time or fuel---corridor-specific speeds and practices such as slow steaming can compress, or even eliminate, NSR's apparent operational advantage over Suez \citep{notteboom2009_fuelcosts,maloni2013_slowsteaming}. Third, although longest, the Cape corridor remains a robust year-round fall-back independent of canals and Arctic seasonality, consistent with the recent diversion of container services around Africa during the Red Sea crisis \citep{Notteboom2024MEL}.

We focus on Europe--Asia trade because it is a major deep-sea liner market and a key arena for debates on the competitiveness of alternative long-haul corridors. Rotterdam--Yokohama spans the Suez, Cape of Good Hope, and NSR corridors and is widely used as a representative origin--destination pair in the NSR--Suez literature. Recent route-comparison studies likewise use Rotterdam--Asia pairs to quantify Arctic versus Suez corridor differences \citep{meza2023_arctic_routes}. 

Macro-waypoints were placed at natural gateway regions (canal entrances, straits, capes). Each successive pair was verified to admit a valid sea-only path in \texttt{searoute}; cases with land intersections or graph failures were iteratively adjusted and then checked visually on interactive charts.

From the reviewed literature, three methodological gaps emerge:

\begin{enumerate}[label=(\roman*),leftmargin=2em]
  \item \textbf{Method isolation is rare.} Comparative studies often remain geometric (GC) or embed routing within complex environmental or economic models, obscuring how routing representation itself biases distance or derived efficiency metrics \citep{liu2010_nsr,schoyen2011_nsrecon,SmithStephenson2013,melia2016_future}.
  \item \textbf{Propagation into emissions is underexplored.} Few works trace routing assumptions through to fuel and CO$_2$ accounting with explicit treatments of main- and auxiliary-engine loads \citep{schroeder2017_emission,chen2021_passenger,karamperidis2022_review,johansson2022_uncertainty,nguyen2023_sensitivity}.
  \item \textbf{Reproducibility is uneven.} Many studies provide limited access to code, parameters, or scenario definitions, hindering cross-study benchmarking and robustness checks.
\end{enumerate}

The primary objective is to quantify how routing representation (GC versus sea-only A* on a 0.5$^\circ$ water mask) alters corridor distances and their translation to voyage time, fuel, and CO$_2$ for Rotterdam--Yokohama across Suez, Cape of Good Hope, and NSR. Secondary objectives are to contrast corridor-typical versus equal-speed baselines, test endpoint sensitivity (for example, Busan versus Yokohama). 

\noindent\textbf{Contributions.} This study provides:
\begin{enumerate}[label=(\alph*),leftmargin=2em]
  \item a corridor-agnostic, reproducible sea-only benchmark for Europe--Asia routing across Suez, Cape of Good Hope, and NSR;
  \item quantified GC$\rightarrow$sea-only deltas under three waypoint philosophies;
  \item transparent propagation to time and fuel/CO$_2$ using explicit main- and auxiliary-engine accounting; and
  \item robustness and endpoint sensitivity analyses identifying when NSR's geometric advantage does---and does not---translate into operational savings.
\end{enumerate}

